\section{R2: Carriers, accumulators, and auxiliary state}\label{sec:r2}

This is the most valuable place to combine a precise mathematical model with
practical synthesis templates. A fold carrier is not just a type: it is a
representation equipped with an initializer, a transition, and an observation
map. Searching the type alone discards most of the useful information.

\question{R2.Q1}{Can finite examples determine a minimal carrier type, or prove that none exists?}
\status{A bounded constructive problem is solvable; the unrestricted request confuses semantic minimal state with a minimal representation in a type grammar.}

For an ordinary list function $h:A^*\to B$, consider a realization
\[
 q([])=z,\quad q(a::xs)=s(a,q(xs)),\quad h(xs)=o(q(xs)),
 \tag{2}\label{eq:fold}
\]
with carrier $K$, initializer $z:K$, step $s:A\times K\to K$, and observer
$o:K\to B$. Define a residual equivalence on lists by
\[
 x\sim_h y\quad\Longleftrightarrow\quad
 \forall u\in A^*,\ h(u\mathbin{+\!+}x)=h(u\mathbin{+\!+}y).
 \tag{3}\label{eq:residual}
\]
The orientation is important: a right fold extends an already-processed tail
by prepending more elements.

\begin{theorem}[The semantic carrier]\label{thm:carrier}
If $q$ realizes $h$ as in \eqref{eq:fold}, then $q(x)=q(y)$ implies
$x\sim_h y$. The quotient $A^*/{\sim_h}$ itself supports a realization of $h$.
Consequently, every reachable realization has at least as many states as this
quotient, in the sense that its reachable states map surjectively onto the
quotient.
\end{theorem}
\begin{proof}
If $q(x)=q(y)$, induction on the prepended list $u$ gives
$q(u\mathbin{+\!+}x)=q(u\mathbin{+\!+}y)$; applying $o$ proves the first claim.
For the quotient, take $z=[\,[]\,]$, $s(a,[x])=[a::x]$, and $o([x])=h(x)$.
The step is well-defined because every test of $a::x$ against $a::y$ is a test
of $x$ against $y$ with context $u\mathbin{+\!+}[a]$. The observer is
well-defined by taking $u=[]$. Finally, map a reachable state $q(x)$ to $[x]$.
The first claim makes this well-defined, and every class is reached.
\end{proof}

This is the residual-state viewpoint behind fold characterizations and
automata minimization, specialized to the proposal's question; it is not a
new universal carrier-inference algorithm \cite{folds}. The quotient may be
infinite, its equivalence may be undecidable, and it need not have a usable
representation in the proposed Lean type grammar. Even a finite semantic
carrier has many representationally equivalent types. Minimal cardinality,
minimal type-expression size, minimal program size, and minimum runtime are
different objectives.

Finite samples do not identify \eqref{eq:residual}. On an infinite domain,
two computable functions can agree on every supplied example but disagree
elsewhere. Without restricting the step and observer, storing the entire
input list and applying $h$ at the end always gives a semantic realization.
Thus an existence claim without a language or cost restriction can be almost
vacuous, while a claim of the intended minimal carrier can be unidentifiable.

A useful decidable version fixes a finite alphabet, a finite carrier of $k$
states, a finite output alphabet, and finitely many observations. Enumerate
transition and observation tables, or encode them as a finite constraint
problem. For a two-symbol alphabet and Boolean outputs, fixing the initial
state to $0$ leaves $k^{2k}2^k$ tables. Searching $k=1,2,\ldots$ gives a smallest
\emph{sample-consistent finite-state machine}. This is constructive, but not
an identification theorem for the intended infinite function.

For the Lean grammar, use bounded type expressions and bounded typed programs
for $z,s,o$, with finite literal and component frontiers. Include proof and
universe bounds when they are searched. Exhaustive enumeration is a decision
procedure only when every membership and observation check in that finite
fragment terminates decisively. A heartbeat timeout must leave the result
unknown, not become a proof that no representation exists.

\paragraph{A sound early refuter.}
Build a graph on observed tails. Connect $x$ and $y$ when some \emph{common}
context $u$ has known unequal outputs for $u\mathbin{+\!+}x$ and
$u\mathbin{+\!+}y$. Every realization assigns different states to adjacent
vertices. A clique of size $k+1$ therefore refutes a carrier with at most $k$
possible reachable states. More strongly, the chromatic number is a lower
bound, although transition congruence can impose further constraints.

Distinct incomplete rows are not enough: a row containing only
$c_0\mapsto\mathsf{false}$ and a row containing only
$c_1\mapsto\mathsf{true}$ are different but compatible. Nor is the set of values
already enumerated for $K$ an upper bound on all possible states. Repeated
application of a tiny transition can build states absent from that enumeration.
Use the actual finite cardinality or a proved reachable-state bound.

\begin{example}[Boolean output, three necessary states]\label{ex:three}
Let $h(xs)$ mean that the length of $xs$ is divisible by three. Tails of lengths
$0,1,2$, tested with contexts of lengths $0,1,2$, give the rows
\[
 (\mathsf{T},\mathsf{F},\mathsf{F}),\qquad
 (\mathsf{F},\mathsf{F},\mathsf{T}),\qquad
 (\mathsf{F},\mathsf{T},\mathsf{F}).
\]
They form a three-clique. A Boolean carrier is impossible even though the
output is Boolean. A three-state carrier exists: increment the length class
cyclically at every step and observe whether it is zero. Its correctness
follows by induction on the list. The finite-model suite additionally
exhausts all one- and two-state machines against 31 binary-list samples of
length at most four; none fits. A three-state witness is found.
\end{example}

\question{R2.Q2}{What replaces an alphabet for polymorphic functions?}
\status{Positions are justified in an audited parametric fragment; equality patterns require an exposed equality operation and a different relational model.}

Suppose a family $h_A:\List A\to\List A$ is natural: it commutes with
\code{map} for every function between element sets. For an input
$xs=[x_0,\ldots,x_{n-1}]$, let $e_n=[0,\ldots,n-1]$ over $\mathsf{Fin}(n)$
and let $v:\mathsf{Fin}(n)\to A$ select the corresponding element of $xs$.
Then naturality gives
\[
 h_A(xs)=\mathsf{map}(v)\bigl(h_{\mathsf{Fin}(n)}(e_n)\bigr).
 \tag{4}\label{eq:naturality}
\]
Thus the output is described by a word of input positions. The same equation
handles repeated elements, because $v$ need not be injective. The empty input
is covered using the empty finite type. This is the precise reason canonical
position-valued tests are valuable, rather than an appeal to polymorphism as
an informal slogan \cite{parametric}.

Equation~\eqref{eq:naturality} is conditional on a naturality theorem. The
surface Lean type alone is not a certificate that arbitrary imported providers,
classical operations, or dictionaries satisfy it. Establish a syntactic
parametric fragment or require relational preservation theorems for admitted
components. A raw \code{BEq A} dictionary in particular is not automatically
lawful or parametric in this sense.

When elements are abstract atoms equipped with equality, the relevant
symmetry is renaming that preserves the available operations. Equality
patterns and a bounded set of stored atoms can then be represented by nominal
or register-style state spaces. Nominal automata supply a mathematical
framework for structured infinite alphabets \cite{nominal}. This does not
make every list transformation finite-state: copying or reversing arbitrary
lists can require unbounded stored data. Separate a finite control state
from symbolic payloads such as saved elements or lists.

Testing canonical inputs up to length $N$ can establish equality on lists of
length at most $N$ under the naturality hypothesis. It does not establish
global equality or equality of all programs of syntactic size at most $d$
unless a separate small-model theorem relates $d$ to $N$.

\question{R2.Q3}{Can fusion run backward to invent auxiliary information?}
\status{Yes as constrained lifting and feature discovery; no general complete inference of an intended helper follows from a few examples.}

AutoLifter is particularly relevant: it decomposes lifting problems into
auxiliary computations and combiners, using a supplied original computation
as a source of observations. Its approximation choices can generate
unrealizable subtasks. Therefore it is neither an unrestricted minimal-type
procedure nor a method requiring only a fixed handful of examples
\cite{autolifter}. Eguchi and colleagues supply another concrete precedent
for inferring auxiliary function types in templates \cite{eguchi}. These are
adaptation starting points rather than reasons to reinvent the entire problem.

For Leant2, represent a proposed carrier together with a symbolic invariant
$J(xs,k)$ explaining the information stored in $k$. Search the three local
obligations
\[
 J([],z),\qquad
 J(xs,k)\Longrightarrow J(a::xs,s(a,k)),\qquad
 J(xs,k)\Longrightarrow C_x(o(k)).
 \tag{5}\label{eq:lift}
\]
Here $C_x$ denotes the required result property at input $xs$, with parameters
made explicit. For a reference implementation it can be equality to the
reference; for a finite contract it is only an observation constraint.
The invariant is a proposal until its obligations are proved.

Use failed symbolic closure to invent features. If the next output requires
an old quantity not reconstructible from the current summary, add that
quantity to the state. For Fibonacci, a proposed state $F_n$ is insufficient
to update uniformly by the elementary recurrence without further information.
The pair $(F_n,F_{n+1})$ admits the closed transition $(a,b)\mapsto(b,a+b)$.
The initializer $(0,1)$ and induction immediately justify the invariant.
This is a useful synthesis template, not a proof that no clever single-natural
encoding exists.

For reversal, the efficient endomorphism carrier has
\[
 K=\List A\to\List A,\qquad z=\mathsf{id},\qquad
 s(a,k)(acc)=k(a::acc),\qquad o(k)=k([]).
\]
An induction gives $q(xs)(acc)=\mathsf{reverse}(xs)\mathbin{+\!+}acc$.
The ordinary list-valued fold also exists but repeatedly appends at the end.
For minimum and maximum together, an optional pair records both the empty
case and two extremes. A total maximum or index operation must make its empty
or out-of-range behavior explicit; inventing a carrier cannot create an
arbitrary element of an uninhabited type.

A complete useful class is obtained by fixing a finite feature vocabulary,
a finite candidate transition language, and decidable closure obligations.
Enumerating feature subsets and their bounded transitions is then complete
for that class. It is an intentionally limited synthesis problem, not a
complete strategy for all catamorphisms over all polynomial functors.

\question{R2.Q4}{How can motives for primitive recursion with parameters be enumerated?}
\status{Enumerate a typed, bounded family of motive templates; completeness is relative to that family and its bound.}

Separate parameters fixed throughout recursion from changing accumulators.
Given a recursor argument $x:I\,\vec\imath$ and a result family $T$, begin with
its goal-directed motive. Extend it in increasing structural cost by a
function space for selected changing parameters, a product for extra outputs,
and an optional state for initialization. Each extension includes its
initializer, observation map, and invariant obligations. This couples motive
selection to the intended computation instead of enumerating arbitrary
functions into a universe.

A useful template has the form
\[
 M(\vec\imath,x)=\Pi (a_1:S_1)\cdots(a_m:S_m),\ K(\vec\imath,x,\vec a),
\]
where $K$ comes from a finite carrier grammar and each binder telescope is
well-formed in dependency order. At a fixed size bound there are finitely
many templates if types, literals, and component choices are also bounded.
Enumerate them once, reject ill-typed ones early, and share residual
constraints between compatible templates. Increasing the bound fairly gives
coverage of the grammar's union, not a terminating procedure for all Lean
inhabitants.

For natural-number recursion, try a single outer recursion capability before
allowing nested recursor use. The proposal attributes an earlier explosion to
nested induction on hypotheses. That is a reason to control capability scope
and template nesting, not to exclude \code{Nat} from the language permanently
\cite{implementation}. Para provides relevant empirical experience with
recursion templates, including function-valued states \cite{para}.

\question{R2.Q5}{Do induction-generalization heuristics transfer to carrier selection?}
\status{Yes as proposal generators, especially when guided by a specific failed step obligation.}

Record why a recursive step failed. If the induction hypothesis is about a
fixed accumulator while the step changes it, propose generalizing that
accumulator. If the step needs two related outputs of the recursive call,
propose tupling. If a symbolic output comparison leaves a recurring expression
unmatched, propose that expression as a feature of the invariant. If a common
subterm appears across failed obligations, propose a helper with that subterm's
type and an example relation.

This is the useful transfer from rippling and lemma speculation: compare the
induction hypothesis with the desired conclusion and isolate a structural
obstruction. It should not be implemented as ``a failed proof proves that this
carrier is wrong.'' Failure may instead reflect insufficient simplification,
a missing library theorem, or inadequate resources. Modern induction work
also investigates generalization inside saturation provers; its search
principles are relevant even when its proof format is not directly reusable
in Lean \cite{inductiongeneralization}.

Use a deterministic, small menu of obstruction-to-template transformations.
Charge every proposed helper and generalized binder in the grammar budget,
and retain the nongeneralized branch. A successful synthesized invariant is
then proved by the local obligations in \eqref{eq:lift}. The highest-priority
experiment is a three-track benchmark: carrier supplied, carrier selected
from a fixed template list, and carrier inferred from failure evidence. This
separates carrier discovery from the cost of synthesizing its transition.
